\documentclass[10pt,oneside]{article}

%%%%%%%%%%%%%
\setlength{\textheight}{8.75in} %Letter is 11in, less 2 for margins, less 0.25 for footer
\setlength{\oddsidemargin}{0.0in} %gets +1inc
\setlength{\evensidemargin}{0.0in} %gets +1inch
\setlength{\textwidth}{6.50in} %Letter is 8.5, less 2 inches for margins
\setlength{\topmargin}{0.5in}
\setlength{\headheight}{0in}
\setlength{\headsep}{0in}
\setlength{\parindent}{0.25in}
%%%%%%%%%%%% 

\usepackage[numbers]{natbib}
\usepackage{graphicx}
\usepackage[colorlinks=true,citecolor=black,urlcolor=blue]{hyperref}

\title{%Hack to get the logo on the PDF front page:
\vspace{-1.5in}
\includegraphics[width=0.5\textwidth]{biopython.jpg}\\
%Hack to get some white space using a blank line:
~\\Biopython Project Update 2014}
\author{
    \underline{Jo\~{a}o Rodrigues}\footnote{Bijvoet Center for Biomolecular Research, Utrecht University, Utrecht, NL. Email: \href{mailto:j.rodrigues@uu.nl}{j.rodrigues@uu.nl}},
    Tiago Antao \footnote{Vector Biology Department, Liverpool School of Tropical Medicine, Pembroke Place, UK},
    Peter Cock\footnote{Information and Computational Sciences, James Hutton Institute (formerly SCRI), Invergowrie, Dundee, UK},
    Eric Talevich\footnote{Department of Dermatology, University of California San Francisco, San Francisco, CA, USA},
    Michiel de Hoon\footnote{Division of Genomic Technologies, RIKEN Center for Life Science Technologies, Yokohama, JP},
		\\
    Wibowo Arindrarto\footnote{Sequencing Analysis Support Core, Leiden University Medical Center, Leiden, NL},
    and the Biopython Contributors}
\date{16\textsuperscript{th} Bioinformatics Open Source Conference (BOSC) 2015, Dublin, Ireland}

\begin{document}
%\pagestyle{empty}
%\thispagestyle{empty}
\maketitle
%\pagestyle{empty}
\thispagestyle{empty}

\vspace{-0.2in}
\noindent
Website: \url{http://biopython.org} \\
Repository: \url{https://github.com/biopython/biopython} \\
License: Biopython License Agreement (MIT style, see \url{http://www.biopython.org/DIST/LICENSE}) \\

The Biopython project is a long-running distributed collaborative effort, 
supported by the Open Bioinformatics Foundation, which develops a freely 
available Python library for biological computation \cite{AppNote}.

We present here details of the recent Biopython release - version 1.65. New 
features include: extended Bio.KEGG and Bio.Graphics modules to support the 
KEGG REST API, as well as parsing, representing, and drawing KGML pathways; 
inclusion of the new NCBI genetic code table 24 (Pterobranchia Mitochondrial) 
and corresponding translation functionality to Bio.Data; improvements to 
Bio.SeqIO (parse and index\_db methods) and Bio.SearchIO (QueryResults objects); 
and refactoring and new functionality of Bio.SeqUtils.MeltingTemp. 
Additionally, we continued our efforts to abide by the PEP8 coding styles, 
namely using lower case module names in new experimental modules.

We are currently preparing a new release - version 1.66 - that will contain 
further work on the Bio.KEGG and Bio.Graphics modules (support for transparency 
in KGML pathways), extended support for the \"abi\" format in Bio.SeqIO, and 
miscellaneous improvements to the test suit and PEP8 coding style. In addition, 
our participation in GSoC 2014 had Evan Parker adding lazy-parsing support for 
Bio.SeqIO, which is currently under review and should be implemented soon.

Finally, complementary to these developments, we created Docker containers and 
added to our GitHub repository. These containers support both Python versions 2 
\& 3 and install all of Biopython's dependencies. They are therefore useful for 
development, but also for teaching due to the inclusion of IPython Notebooks.

\begin{thebibliography}{}

\bibitem[Cock {\it et al}., 2009]{AppNote}Cock, P.J.A., Antao, T., Chang, J.T., Chapman, B.A., Cox, C.J., Dalke, A., Friedberg, I., Hamelryck, T., Kauff, F., Wilczynski, B., de Hoon, M.J. (2009) Biopython: freely available Python tools for computational molecular biology and bioinformatics. {\it Bioinformatics} {\bf 25}(11) 1422-3. \href{http://dx.doi.org/10.1093/bioinformatics/btp163}{doi:10.1093/bioinformatics/btp163}

\end{thebibliography}

\end{document}
